% from ega4-2-fr.tex
\oldpage[IV-2]{201}

au-dessus de l'idéal maximal de $A$, et $B=C_{\mathfrak n}$, il s'agit de
montrer que pour tout idéal premier $\mathfrak q$ de $\widehat B$, dont
l'image réciproque dans $B$ est 0, l'anneau
$\widehat B_{\mathfrak q}$ est \emph{régulier}. On peut d'ailleurs remplacer
$A$ par son image dans $C$, qui est un anneau local complet (comme quotient
de $A$) et intègre. La conclusion à prouver est alors conséquence du lemme
suivant, plus général en apparence.
\end{proof}

\begin{lemma}[7.4.4.2]
\label{IV.2.s7.4.4.2-fr}
Soient $A$ un anneau local noethérien intègre et complet, $k$ son corps
résiduel, $C$ un anneau intègre contenant $A$, tel qu'il existe $t\in C$
pour lequel $C$ soit une $A[t]$-algèbre finie. Soit $\mathfrak n$ un idéal
maximal de $C$ au-dessus de l'idéal maximal $\mathfrak m$ de $A$; posons
$B=C_{\mathfrak n}$, $X=\operatorname{Spec}(B)$,
$B'=\widehat B$, $X'=\operatorname{Spec}(\widehat B)$; si
$U=\operatorname{Reg}(X)$, $U'=\operatorname{Reg}(X')$, et si
$f:X'\to X$ est le morphisme canonique, on a alors
$f^{-1}(U)\subset U'$.
\end{lemma}

\begin{proof}
L'assertion à démontrer pour obtenir (7.4.4.1) se déduira de ce lemme en
observant que, puisque $B$ est intègre, le point générique de $X$ appartient
à $U$.

On notera que puisque $C$ est une $A$-algèbre de type fini, et $\mathfrak n$
un idéal maximal de $C$, le corps résiduel $k'$ de
$C_{\mathfrak n}=B$ (donc aussi de $B'$) est une extension finie de $k$
(I, 6.4.11 et 6.4.2).

Si l'on pose $Y=\operatorname{Spec}(C)$, il résulte de (6.12.8) que
$\operatorname{Reg}(Y)$ est ouvert dans $Y$; comme les anneaux locaux de
$X$ sont des anneaux locaux de $Y$ (I, 2.4.2), on a
$U=X\cap\operatorname{Reg}(Y)$, donc $U$ est ouvert dans $X$; d'autre part
(6.12.7) $U'$ est ouvert dans $X'$, donc $S'=X'-U'$ est fermé; par suite
$S'\cap f^{-1}(U)$ est localement fermé dans $X'$, et tout revient à
prouver que cet ensemble est vide. On sait (5.1.10) que dans le cas
contraire, il existerait dans $S'\cap f^{-1}(U)$ un idéal premier
$\mathfrak p'$ tel que $\dim(B'/\mathfrak p')\leqslant1$. Remarquons d'abord
que $\mathfrak p'$ ne peut être l'idéal maximal $\mathfrak mB'$ de $B'$,
où $\mathfrak m$ est l'idéal maximal de $B$. En effet, cela signifierait que
$B=B_{\mathfrak m}$ serait régulier, donc aussi $B'=\widehat B$ (0,
17.1.5), et on aurait par suite $\mathfrak mB'\in U'$ contrairement à
l'hypothèse. On doit donc avoir $\dim(B'/\mathfrak p')=1$. Posons
$\mathfrak p=B\cap\mathfrak p'$; par hypothèse $B_{\mathfrak p}$ est
régulier, mais $B'_{\mathfrak p'}$ ne l'est pas; comme
$B'_{\mathfrak p'}$ est un $B_{\mathfrak p}$-module plat, il résulte de
(6.5.2) que la fibre $Z$ de $f$ au point $\mathfrak p$ n'est pas régulière
au point $\mathfrak p'$. Montrons qu'on peut se ramener au cas où
$\mathfrak p=0$. En effet, dans le cas général, si l'on pose
$\mathfrak q=\mathfrak p\cap C$, $\mathfrak r=\mathfrak p\cap A$,
$C/\mathfrak q$ est une $(A/\mathfrak r)[\bar t]$-algèbre finie (où
$\bar t$ est la classe de $t$ mod. $\mathfrak q$); comme
$\mathfrak p=\mathfrak qB$, $B/\mathfrak p$ est égal à
$(C/\mathfrak q)_{\mathfrak n/\mathfrak q}$, et
$\mathfrak n/\mathfrak q$ est un idéal maximal de $C/\mathfrak q$
au-dessus de l'idéal maximal $\mathfrak m/\mathfrak r$ de
$A/\mathfrak r$; on voit donc que les hypothèses de (7.4.4.2) sont encore
remplies par $A/\mathfrak r$, $C/\mathfrak q$ et $B/\mathfrak p$, et comme
le complété de $B/\mathfrak p$ est $B'/\mathfrak pB'$, cela prouve notre
assertion. Supposons donc $\mathfrak p=0$, de sorte que $Z$ est la fibre
générique et que l'homomorphisme $B\to B'/\mathfrak p'$ est injectif.
Posons $V=B'/\mathfrak p'$, et distinguons deux cas :

\emph{I) $V$ est une $A$-algèbre finie.} --- Comme $B\subset V$, $B$ est
\emph{a fortiori} une $A$-algèbre finie, et comme $A$ est complet, il en est
de même de $B$ (Bourbaki, \emph{Alg. comm.}, chap. IV, \S~2,
n\textsuperscript{o}~5, cor. 3 de la prop. 9), d'où $B'=B$,
$\mathfrak p'=0$, donc $B'_{\mathfrak p'}$ est un corps, et par suite un
anneau régulier, contrairement à l'hypothèse.

\emph{II) $V$ n'est pas une $A$-algèbre finie.} --- Comme l'anneau local
$A$ est complet, cela implique que $V$ n'est pas une $A$-algèbre
quasi-finie (0$_{\mathrm I}$, 7.4.1 et 7.4.2); mais par hypothèse le corps
résiduel $k'$ de $V$ est extension finie du corps résiduel $k$ de $A$,

